\noindent\textbf{(a)} Es un factor integrante en $xy(x+y)\ne0$, para todos $a,b$. Un potencial es $F=b\ln|xy|+(a-b)\ln|x+y|$.

\noindent\textbf{(b)} Si $b\ne0$, vale $\mu_2=|s|^{a/b-2}$ en $s\ne0$. Si $b=0$ y $a\ne0$, no existe tal factor en una región abierta. Si $a=b=0$, cualquier función no nula adecuada sirve.

\noindent\textbf{(c)} Si $a=b=0$, toda función $C^1$ sirve. Si $b=0$, $a\ne0$: $y=0$ o $y=C-x$. Si $b\ne0$, $a=-b$: $y=Cx$. En los demás casos, fuera de $xy(x+y)=0$, $b\ln|xy|+(a-b)\ln|x+y|=C$, además de $y=0$ y $y=-x$; se toman ramas gráficas de la EDO original. En este último caso ($b\ne0$, $a+b\ne0$), para atravesar $x=0$ se exige $y(0)=0$: se pueden unir ramas de la lista anterior cuyos valores tiendan a cero y cuyas derivadas tiendan a una misma pendiente, $0$ o $-1$. Estas condiciones son necesarias y suficientes; las constantes a izquierda y derecha pueden ser distintas.
